\section{Discussion}
\label{sec:discussion}

\subsection{Behavioral Mastery vs.\ Representational Structure}
\label{sec:disc_mastery}

The most striking finding is the disconnect between behavioral performance and internal representation.
\gptmodel achieves 100\% accuracy on all French number tasks, including vigesimal numbers---a result that refutes the initial hypothesis of 10--30\% lower accuracy on vigesimal forms.
This perfect performance likely reflects the sheer volume of French text in modern training corpora: the model has memorized the complete mapping between vigesimal words and their values.

Yet \mistral's internal representations reveal that this memorization preserves the linguistic structure of the vigesimal system.
Numbers 80--99, which all share the \quatrevingt prefix, cluster together in representation space regardless of their numeric distance.
The model ``thinks in vigesimal''---organizing numbers by their linguistic base---even though it ``answers in decimal'' with perfect accuracy.
This finding aligns with \citet{marjieh2025number}'s observation that LLM number representations entangle string-level and numerical properties, and extends it to a language where the string-number mapping is non-trivially compositional.

\subsection{The Belgian French Advantage}
\label{sec:disc_belgian}

Belgian French serves as a natural ablation of the vigesimal system.
By replacing \emph{soixante-dix} (60+10) with \emph{septante} (70) and \emph{quatre-vingt-dix} (4$\times$20+10) with \emph{nonante} (90), Belgian French removes the implicit arithmetic while preserving the language.
We find that Belgian French produces representations with the highest ordinal correlation ($\rho = 0.591$ vs.\ $0.520$ for standard French), confirming that \textbf{linguistic regularity directly improves the ordinal structure of numeric representations}.

The near-identical per-number cosine similarity between French and Belgian French (0.995) is equally notable.
The model maps both \emph{quatre-vingt-treize} and \emph{nonante-trois} to essentially the same point in representation space, suggesting it recognizes them as the same number.
However, the pairwise distance structure differs ($\RSA$ correlation = 0.890, not 1.0), because the vigesimal prefixes create distortions in the French RDM that are absent in the Belgian French RDM.

\subsection{Why French Words Beat Digits in Probing}
\label{sec:disc_words}

The counterintuitive finding that French words are more linearly decodable ($R^2 = 0.979$) than digit strings ($R^2 = 0.943$) deserves explanation.
We hypothesize that the prompt template ``\emph{Le nombre est [X]}'' activates a semantic frame that enriches the representation with numerical meaning.
When the model processes \emph{Le nombre est quarante-deux}, it has already been primed to interpret the following tokens as a number, producing a representation that is optimized for numeric content.
In contrast, a bare digit string like ``42'' in the same template may be more ambiguous---it could be a code, an identifier, or a quantity---leading to a less numerically structured representation.

This result has practical implications: \textbf{linguistic context can improve rather than degrade numeric representations}, contrary to the assumption that words introduce noise relative to digits.

\subsection{Limitations}
\label{sec:limitations}

\para{Single model for representational analysis.}
We analyze only \mistral for internal representations.
The linguistic clustering effect may differ in models with different architectures, tokenizers, or training data distributions.
Extending to Llama~3, Qwen, or other multilingual models would strengthen these findings.

\para{Last-token extraction.}
We extract hidden states at the last token position only.
Alternative aggregation strategies (mean pooling, attention-weighted pooling) might yield different representational structures, particularly for multi-token vigesimal words where information is distributed across tokens.

\para{Context sensitivity.}
Our results are specific to the prompt template ``\emph{Le nombre est [X]}.''
Other contexts---arithmetic problems, counting sequences, free-form text---may produce different representational patterns.

\para{Limited number range.}
We focus on 0--99.
The vigesimal system also applies to hundreds (\eg \emph{deux-cent-quatre-vingt-dix-sept} = 297), where the compositional complexity increases further.

\para{No causal intervention.}
We observe representational differences but do not establish causality.
Activation patching \citep{levy2024language} could test whether the linguistic clustering causally affects downstream outputs, but we leave this to future work.
